\documentclass[]{article}

\usepackage{amsmath}
\usepackage{multirow}
\usepackage{natbib}
\usepackage{graphicx} 


\begin{document}

The discovery of neutrino oscillations has unequivocally demonstrated that neutrinos have mass, providing the first definitive evidence of physics beyond the Standard Model. While oscillation experiments have measured the two mass-squared differences with remarkable precision, they are insensitive to the absolute neutrino mass scale and the ordering of the mass eigenstates. This ordering, known as the neutrino mass hierarchy, remains one of the most significant open questions in particle physics. The hierarchy can be Normal (NH), where two lighter masses are separated from a heavier one ($m_1 < m_2 \ll m_3$), or Inverted (IH), with one light mass separated from two heavier, nearly degenerate ones ($m_3 \ll m_1 < m_2$). Resolving this ambiguity is crucial for developing theories of fermion mass generation and has profound consequences for experimental searches for new physics in the lepton sector.The determination of the mass hierarchy is particularly critical for the ongoing search for neutrinoless double-beta decay ($0\nu\beta\beta$). The observation of this lepton-number-violating process would prove that neutrinos are their own antiparticles (Majorana fermions). The rate of this decay is governed by the effective Majorana mass, $m_{\beta\beta}$, an observable that depends on the absolute neutrino masses, mixing parameters, and two unknown CP-violating phases. The two hierarchies predict starkly different ranges for $m_{\beta\beta}$. The Inverted Hierarchy sets a lower limit on $m_{\beta\beta}$ that falls squarely within the sensitivity of next-generation experiments. In contrast, the Normal Hierarchy permits values of $m_{\beta\beta}$ that could be vanishingly small, potentially placing the detection of $0\nu\beta\beta$ beyond the reach of any foreseeable technology. A definitive measurement of the hierarchy is therefore essential for interpreting the results of these multi-billion-dollar experimental programs and for guiding future research strategies.Cosmology offers a powerful and independent method to address this question. Massive neutrinos behave as a form of hot dark matter, suppressing the growth of large-scale structures on scales smaller than their free-streaming length. Cosmological observables, such as the galaxy power spectrum and the cosmic microwave background, are therefore sensitive to the sum of the neutrino masses, $\Sigma m_\nu$. While oscillation experiments do not determine this sum, they impose a strict lower bound on it for each hierarchy. Using current oscillation data, the minimum allowed value is $\Sigma m_\nu \ge 0.06\,\text{eV}$ for the Normal Hierarchy, whereas for the Inverted Hierarchy, it is $\Sigma m_\nu \ge 0.1\,\text{eV}$. This difference provides a clear and testable prediction: if cosmological data can robustly constrain the total neutrino mass to be well below $0.1\,\text{eV}$, the Inverted Hierarchy would be decisively disfavored.In this work, we perform a rigorous Bayesian model comparison between the Normal and Inverted Hierarchies using the latest cosmological data from the Dark Energy Spectroscopic Instrument (DESI) Data Release 2. The unprecedented precision of DESI provides the tightest cosmological constraints on $\Sigma m_\nu$ to date, enabling a definitive test of the hierarchy for the first time. By combining the cosmological likelihood from DESI with the most recent global analysis of neutrino oscillation data, we compute the Bayesian evidence for each hierarchy. We meticulously assess the robustness of our conclusions by employing distinct, well-motivated prior probability distributions for the unknown neutrino mass parameters. Furthermore, we quantify the sensitivity of our results to the underlying cosmological model by extending our analysis beyond the standard flat $\Lambda$CDM paradigm. Our analysis provides a statistically robust and quantitative answer to the question of the neutrino mass hierarchy, yielding direct and impactful consequences for the future of neutrinoless double-beta decay experiments and our broader understanding of fundamental physics.

\end{document}
                